\section{Related Work}
\label{sec:related-work}

\subsection{Serving Architectures and Scheduling Mechanisms}
\label{sec:related-serving}

A substantial body of systems work has improved LLM inference throughput and
latency through scheduling and memory-management mechanisms. Orca introduced
iteration-level scheduling at the granularity of individual decoding steps
\citep{yu2022orca}. vLLM introduced PagedAttention, a paged
KV-cache allocator that raised achievable batch sizes and became a widely
adopted serving-system reference point \citep{kwon2023vllm}. Sarathi-Serve
interleaves chunked-prefill computation with ongoing decode steps to reduce
head-of-line blocking under continuous batching \citep{agrawal2024sarathi}.
Splitwise and DistServe disaggregate the compute-bound prefill phase and the
memory-bound decode phase onto separate hardware pools
\citep{patel2024splitwise,zhong2024distserve}; Mooncake extends this
direction with a KV-cache-centric disaggregated architecture that schedules
around a shared, elastic KV cache \citep{qin2025mooncake}. Llumnix
adds dynamic cross-instance request migration and rebalancing
\citep{sun2024llumnix}, and DynamoLLM reconfigures an inference cluster's
resources for energy and cost efficiency under latency SLOs
\citep{stojkovic2025dynamollm}.
For geographically distributed LLM inference, \citet{sun2025geodistributed}
study block placement and request routing over heterogeneous GPU servers,
combining experimentally validated performance models with optimization
and simulation. That work is directly adjacent to LSSP in treating LLM
inference as a resource-allocation performance problem, but its object is
placement and routing optimization for one distributed serving
architecture rather than the portability of scheduler rankings across
independent workloads and metrics.

A second cluster of work targets scheduling \emph{policy} design more
specifically. VTC (Virtual Token Counter) defines an explicit token-level
fairness objective for continuous-batching schedulers with a provable
service-difference bound \citep{sheng2024vtc}. $S^3$
predicts each request's output length to pack KV-cache memory more tightly
at admission time \citep{jin2023s3}, whereas
\citet{efficientllmscheduling2024} learns a pairwise ranking model over
waiting requests to prioritize decode order. Recent systems diversify the
mechanism space further along four lines: preemptive, admission-control,
and SLO-aware scheduling under imprecise length information or
heterogeneous per-request SLOs, as in FastServe, JITServe, and SCORPIO
\citep{fastserve2024,jitserve2025,scorpio2025}; learned-ranking,
multi-objective, and single-score policies, as in PARS, the SLAI/RAD family,
and LMetric \citep{pars2025,slai2025,lmetric2026}; request partitioning,
hybrid KV-cache/hidden-state representations, and multi-priority
engine/service-layer scheduling, as in Libra, Apt-Serve, and ProServe
\citep{libra2026,aptserve2025,proserve2025}; and tail-latency,
multi-tenant, and agentic scheduling, including prediction-free
interpolation between first-come-first-served and shortest-job-first with
cache-aware preemption \citep{li2026beyond},
hierarchical multi-agent scheduling under non-stationary multi-tenant load
\citep{liu2026hmas}, and session-centric routing that trades load balance
against cache locality \citep{wang2026smetric}. Each of these
papers reports an improvement over one or more baselines on
the workload(s) chosen for its own experiments; none asks whether the
\emph{relative ordering} it reports would
be preserved under an independently sourced workload, a different metric,
or a different operating load -- the question this benchmark is built
around (\S\ref{sec:related-position}).

Classical and data-center scheduling work offers useful context for the
distinction between designing a scheduling policy and evaluating the
robustness of scheduler conclusions. For example,
\citet{chen2026quickswap} design and analyze non-preemptive policies for
multiserver jobs and evaluate them on data-center traces, addressing the
mean-response-time consequences of high-utilization scheduling decisions.
LSSP does not propose a new scheduler in that line; it asks whether
comparative scheduler-performance conclusions remain stable when the
workload source, operating region, or metric changes.

\subsection{Production Workloads and Characterization}
\label{sec:related-workloads}

A parallel literature documents that LLM-serving workloads observed in
production are themselves heterogeneous. BurstGPT releases 10.31 million
traces collected from a regional deployment of Microsoft's Azure OpenAI GPT
service over 213 days, characterizing request burstiness, conversation
structure, and response-length behavior, and is now published as a full
KDD~2025 paper rather than only as a preprint \citep{burstgpt2025}. Two
official Microsoft releases document Azure's own internal LLM-inference
traffic: Splitwise analyzes production conversation and code-completion
traces collected on Azure in November 2023 \citep{patel2024splitwise}, and
DynamoLLM analyzes a longer, one-week Azure trace collected in May 2024
\citep{stojkovic2025dynamollm}. \citet{kvcache_wild2025} characterizes
KV-cache reuse patterns at a large cloud provider, finding that reuse
probability and lifespan are workload-category-dependent in ways that
synthetic workload studies had not captured. \citet{ayearinllmserving2026}
reports a one-year longitudinal characterization of production traffic
spanning workload evolution, caching behavior, and load balancing (a recent
preprint at the time of writing; \S\ref{sec:related-position}).
ServeGen instead characterizes a worldwide production inference service in
order to build a principled per-client workload \emph{generator}, motivated
explicitly by the position that no single fixed trace is representative
enough on its own \citep{servegen2026}; because it characterizes the same
underlying Alibaba Cloud platform as this project's Bailian/Qwen source, we
do not count it as an independent workload provider (\S\ref{sec:workloads}).
TraceLab characterizes a distinct workload family entirely -- coding-agent
sessions rather than single-turn or short-multi-turn chat requests -- and is
used here only as out-of-distribution context, not as one of this
project's three primary sources.

\subsection{Simulators and Benchmark Frameworks}
\label{sec:related-simulators}

Vidur is the closest existing infrastructure precedent to LSSP: a
large-scale simulation framework purpose-built for exploring LLM-inference
system configurations at a fidelity and speed unavailable from
running full experiments on real hardware \citep{vidur2024}. LLMServingSim
and its successor LLMServingSim~2.0 provide a related but architecturally
different hardware/software co-simulation infrastructure, adding
iteration-granularity simulation with algorithmic-redundancy reuse and,
in the 2.0 revision, explicit support for heterogeneous accelerators and
disaggregated serving techniques \citep{llmservingsim2024,llmservingsim2_2025}.
Two recent 2026 benchmarking frameworks expand this domain: CCL-Bench 1.0
\citep{ding2026cclbench} introduces a diagnostic, trace-based evaluation
of LLM training and serving infrastructure to record fine-grained execution
evidence of kernels and communication events, while XPerf \citep{wang2026xperf}
benchmarks serving systems under agentic workloads by employing fine-grained
trace replay and full-stack profiling to handle agentic nondeterminism.
These simulators and benchmarking frameworks are built to let researchers
explore serving-system design points cheaply or load-test systems under
agentic dependencies; none of them is built around, or reports, a systematic
benchmark of \emph{comparative scheduler ranking portability} across
independently sourced workloads, operating regions, and metrics in the sense
this paper defines. LSSP's own simulator (\S\ref{sec:methodology}) is
project-local infrastructure, not a contribution positioned against this
literature; we make no claim that its output is equivalent to production
hardware latency or throughput, only that it supports a fixed, controlled
benchmark protocol for relative comparisons (\S\ref{sec:methodology}).

\subsection{Workload Dependence and Benchmark Sensitivity}
\label{sec:related-dependence}

Prior work has already established, individually, that LLM-serving
workloads are heterogeneous (\S\ref{sec:related-workloads}) and that
individual serving mechanisms and schedulers are sensitive to batching
structure, prefill/decode balance, load level, and service-level
objectives (\S\ref{sec:related-serving}). We treat both of these as accepted
background, not as a novel contribution of this work. Closest to this
paper's motivation, \citet{papaioannou2024workload} show that a serving
system's measured performance and the relative standing of design choices
can themselves shift depending on which workload is used to evaluate them,
arguing for workload-choice awareness as a first-class methodological
concern; that result is stated for single-system evaluation. In classical
data-center dispatching, \citet{yildiz2026dispatching} compare Round Robin,
Join-Idle-Queue, and Least-Work-Left on Google ClusterData v3 and Alibaba
Cluster Trace v2018 using analytical models and trace-driven simulation.
Their results show that production workload structure and the job-level
versus task-level trace representation can change dispatching-policy
behavior and, in some production-trace cases, the policy ordering. This is
the closest non-LLM precedent for the evaluation concern LSSP studies. The
distinction is that Yildiz et al.\ diagnose model--trace discrepancies and
dispatcher behavior for data-center workloads, whereas LSSP makes the
portability of a \emph{comparative ordering} among a fixed panel of
schedulers the primary benchmark object, asking whether that ordering is
preserved across independently sourced workloads, operating regions, and
metrics in LLM serving -- the distinction this paper's Introduction
develops between single-system sensitivity and comparative-ranking
portability.

\subsection{Ranking Portability and Benchmark Methodology}
\label{sec:related-methodology}

The statistical methodology this benchmark uses -- Kendall's tau-b and
Spearman's rho for pairwise ranking agreement, paired comparison and
block-bootstrap confidence intervals for ranking-uncertainty
quantification, and a probability-of-correct-selection framing for
benchmark sample complexity -- draws on standard nonparametric statistics
and comparative-algorithm-ranking methodology (\S\ref{sec:methodology});
these are general statistical tools, not literature claims specific to LLM
serving, and are cited here only as the methods this benchmark's analysis
plan applies to its own, independently generated data. The closest
methodological precedent we identified is from the classical parallel-job-
scheduling literature: \citet{zakay2014resampling} partition a real
scheduler workload trace into independent per-user subtraces and recombine
them to resample varied-but-realistic workloads, evaluating how a
\emph{single} parallel job scheduler's measured performance responds to
workload resampling. LSSP's sample-complexity question
(\S\ref{sec:sample-complexity}) is the same evaluation instinct --
using resampling to ask how much workload evidence a scheduler-performance
conclusion actually rests on -- applied to a different object: the
portability of a \emph{comparative} ranking among multiple schedulers
across independently sourced workloads, rather than one scheduler's
sensitivity to resampled variants of a single trace.

\subsection{Position of LSSP}
\label{sec:related-position}

Prior work generally studies a scheduler, a serving system, a
simulation environment, or a workload generator/characterization in its own
right. LSSP instead treats the portability of \emph{comparative} scheduler
rankings as the primary object of study across independently sourced
LLM-serving workload sources, operating regimes, and metrics, using a
fixed, mechanism-diverse scheduler panel and a preregistered protocol, while
explicitly quantifying ranking uncertainty and benchmark sample complexity.
A primary-source literature search targeting exactly this
question (scheduler-ranking stability/portability, cross-workload
scheduler evaluation, benchmark sample complexity and representativeness
for LLM-serving systems, current through September 2026) did not surface
prior work with this primary object. The closest precedents differ in
object: \citet{yildiz2026dispatching} show on Google and Alibaba
data-center traces that workload structure and trace granularity can affect
dispatcher behavior and policy ordering, but do not study LLM-serving
schedulers, cross-metric ranking portability, supported reversal analysis,
or benchmark sample complexity; \citet{zakay2014resampling}
(\S\ref{sec:related-methodology}) evaluate one scheduler's
sensitivity to workload resampling, not comparative ranking across
schedulers; Vidur and LLMServingSim are simulation infrastructure
(\S\ref{sec:related-simulators}); and the production-characterization
literature (\S\ref{sec:related-workloads}) establishes the workload
heterogeneity this design responds to without studying cross-source
ranking portability. This wording is intentionally narrower than an
unconditional claim of priority: we do not claim to be the first to study
LLM-serving workload dependence or scheduler sensitivity, only that we did
not identify a prior benchmark treating cross-source, cross-metric,
cross-load-region scheduler-ranking portability as its primary object of
study.
